% Generated draft: requires mathematical and source review.
\documentclass[11pt]{article}
\usepackage[margin=25mm]{geometry}
\usepackage{fontspec}
\setmainfont{Latin Modern Roman}
\usepackage{amsmath,amssymb,mathtools,unicode-math}
\setmathfont{Latin Modern Math}
\usepackage{enumitem,needspace,xurl,hyperref}
\setlength{\parindent}{0pt}
\setlength{\parskip}{5pt}
\setlength{\emergencystretch}{3em}
\newcommand{\N}{\mathbb{N}}
\newcommand{\Z}{\mathbb{Z}}
\newcommand{\Q}{\mathbb{Q}}
\newcommand{\R}{\mathbb{R}}
\newcommand{\C}{\mathbb{C}}
\newcommand{\F}{\mathbb{F}}
\newcommand{\A}{\mathbb{A}}
\newcommand{\PP}{\mathbb{P}}
\newcommand{\E}{\mathbb{E}}
\newcommand{\eps}{\varepsilon}
\newcommand{\dd}{\,\mathrm{d}}
\begin{document}
\section*{UnsolvedMath ID: 30002897}
\textbf{Source problem code:} catalogue entry
\subsection{Definitions and mathematical statement}
For a von Neumann algebra \(M\) and a subset \(S\subseteq M\), define the von Neumann algebra generated by \(S\) inside \(M\), using the unital convention, by
\[
W_M^*(S)=\bigcap\{N\subseteq M:N\text{ is a von Neumann subalgebra},\ 1_N=1_M,\ S\subseteq N\}.
\]
The supplied question is: is the tail algebra \(T_\phi\) generated as a von Neumann algebra by \(D_\phi\)? Provided that the intended objects are subalgebras of a common ambient von Neumann algebra \(M\), with the intended common-unit convention, this asks whether
\[
T_\phi=W_M^*(D_\phi).
\]
The supplied record does not define \(\phi\), its admissible class, the underlying algebra or representation, the tail algebra \(T_\phi\), or \(D_\phi\). It also does not state the quantifier over \(\phi\). Consequently, the displayed equality is a schematic rendering of the question, not a complete quantified formulation. The title refers to diagonal subalgebras, but supplies no definition identifying \(D_\phi\). These missing definitions and conventions must be recovered from the cited source before a precise standalone statement can be finalized.
\subsection{Short English statement}
The question asks whether the specified subalgebra \(D_\phi\) generates the entire tail von Neumann algebra \(T_\phi\). The supplied record leaves the objects, hypotheses, and quantification over \(\phi\) undefined, so the formulation remains provisional.
\subsection{Sources}
\begin{itemize}
\item Local UnsolvedMath registry, supplied record for permanent ID \(30002897\), source problem code \texttt{catalogue entry}, titled ``Generation of Tail Algebras by Diagonal Subalgebras,'' in the Oberwolfach Reports Open Problems collection. No individual proposer is supplied; the recorded proposed year is 2015.
\item Supplied source citation: \emph{Free Probability Theory (2016)}, Oberwolfach Reports, DOI 10.4171/owr/2015/28, pp.\ 1571--1629. \url{https://doi.org/10.4171/owr/2015/28}. The supplied background describes these reports as records of mathematical workshops and their open-problem discussions.
\item Supplied literature source: ``Oberwolfach report 28.'' \url{https://doi.org/10.4171/owr/2015/28}.
\item Supplied literature source: ``Quantum symmetric states on free product C-star algebras.'' \url{https://arxiv.org/abs/1305.7293}.
\item Literature assessment dated exactly \texttt{2026-08-21}: the supplied curation labels the problem open and states that the later report still lists generation of the von Neumann tail algebra \(T\) by \(D\) as open, and that related de Finetti and amalgamated-freeness theorems do not establish this generation statement. This is supplied dated curation, not an independently verified literature assessment. The supplied background explicitly calls for independent verification of time-sensitive status and cited claims.
\end{itemize}
\end{document}
